\section{More on ATP usage}\label{atp-app}

\subsection{Basic usage}
When running an ATP or finite model builder on a nontrivial problem, the two key recommendations are: 1) Conduct several runs with different search-parameter settings, and 2) provide enough computational resources for each search: memory\footnote{If a memory limit is not specified, \emph{Vampire} will try to use as much RAM as possible, but by default \emph{Prover9}-\emph{Mace4} sets memory limits which are rather low for current standards, so we advise raising them before starting a long computation.}, number of processing cores\footnote{\emph{Vampire} can make use of as many cores as the user specifies. To take advantage of multicore processors with \emph{Prover9}-\emph{Mace4}, one can run several terminal or GUI instances. Moreover, each GUI window allows to run \emph{Prover9} and \emph{Mace4} simultaneously on the same problem, and they run on different cores.}, number of user instructions executed, and especially, time. Along the ETP, some proofs or models that could not be found in under 1 second with a given configuration, could be obtained in under (say) 8 seconds without altering the configuration; other implications initially required runs lasting several minutes, or even hours. Once a proof is found by some configuration, a short and quick proof can typically be found by tweaking that configuration.

The configuration of the search parameters is a difficult art that has become more sophisticated as ATPs have grown in complexity. To address difficult problems, we strongly recommend acquiring a solid understanding of the available configuration options and their effects, enabling the search space to be tailored as closely as possible to the problem at hand. Currently, options allow control over numerous aspects of each proving stage, including search limits, preprocessing steps, inference rules, formulas ordering and weighting, etc.

\subsection{Flow control}\label{flow control}
\emph{Vampire} is equipped with a user-friendly and powerful standard mode, the CASC mode (or CASC SAT mode for finite model building) that in many situations avoids the need of configuring specific search options for the given problem. This mode invokes a sequence of strategies (different option configurations) with assigned time limits. The time given to each strategy is important: counterintuitively, giving less time to the CASC mode may end up producing a faster proof: there are proofs that can be found when the time limit is set to less than $5\%$ of the time \emph{Vampire} needs to find a proof without the time limit~\cite{DBLP:conf/cav/KovacsV13}. At least two factors contribute to this effect: the correct strategy is explored sooner because less time is spent on each strategy, and the behaviour of \emph{Vampire}'s default saturation algorithm, the limited resource strategy (LRS) algorithm.

On the other hand, \emph{Prover9} itself does not have the capability of running several schedules in a row, although it lets the user implement some rudimentary strategies through commands called \emph{actions}. Also, a separate program called \emph{FOF-Prover9} includes a preprocessing step that attempts to reduce the problem to independent subproblems, possibly reducing the overall time significantly.

\subsection{The weight-sos limit strategy}\label{weight_strategy}

In \emph{Prover9}, each clause is assigned a weight depending on its length (and other customizable parameters), with newly generated clauses exceeding the \verb|max_weight| limit being discarded. In addition, the size of the list of clauses awaiting processing (the SOS list) is controlled with the \verb|sos_limit| parameter, with most newly generated clauses being discarded once the list is full. Consequently, the size and shape of the search space can be partially controlled through these parameters. Ideally, we would want to use the smallest values of \verb|max_weight| and \verb|sos_limit| that still guarantee a proof to be contained in the search space. Thus smaller values are preferable, provided they are not so low that the search is exhausted without finding a proof.

Since there are more parameters affecting the search, and different proofs may be reachable depending on the configuration, the system's behaviour with respect to \verb|max_weight| and  \verb|sos_limit| is not straightforward, with several casuistics arising. When \verb|sos_limit| is fixed, the proof-finding time tends to rise with \verb|max_weight| until reaching a plateau, with the length of this increasing phase extending for higher \verb|sos_limit| values (see \Cref{figure:weight_time}). For this reason we recommend lowering \verb|sos_limit| from its default value of 20000 to substantially smaller values.

\begin{figure}
\centering
\includegraphics[width=225pt]{450_max_weight}
\includegraphics[width=230pt]{102744082_max_weight}
\caption{Proof-finding time as a function of \texttt{max\_weight}, for several values of \texttt{sos\_limit}.}
\label{figure:weight_time}
\end{figure}

On the other hand, when \verb|max_weight| is fixed, the smallest \verb|sos_limit| values that still yield proofs typically produce a chaotic transient phase with higher proof-finding times, after which a better-defined relation between proof-finding time and \verb|sos_limit| emerges, which tends to follow either a near-plateau pattern (\Cref{figure:sos_limit_time}a)), or a monotonically increasing trend (\Cref{figure:sos_limit_time}b)). We recommend avoiding an excessively low \verb|sos_limit|, in order to prevent the onset of the transient phase.

\begin{figure}
\centering
\includegraphics[width=305pt]{450_sos_limit_time}
\includegraphics[width=310pt]{650_sos_limit_time}
\caption{Proof-finding time as a function of \texttt{sos\_limit}, for different values of \texttt{max\_weight}. a) Stabilization as plateau, same for different weights. b) Different behaviours for different weights, with one trend monotonically increasing.}
\label{figure:sos_limit_time}
\end{figure}

Other relevant characteristics related to the complexity of the resulting proof (such as proof level, length, weight, etc.\@) may also vary in nontrivial ways depending on the chosen \verb|sos_limit| value (see \Cref{figure:sos_limit_proof_data}).

\begin{figure}
\centering
\includegraphics[width=230pt]{450_sos_limit_proof_data}
\includegraphics[width=230pt]{650_sos_limit_proof_data}
\caption{Proof-complexity indicators as a function of \texttt{sos\_limit}.}
\label{figure:sos_limit_proof_data}
\end{figure}

Remarkably, \emph{Prover9} can establish all positive implications between laws of order up to 4 using \verb|max_weight| 55 and \verb|sos_limit| 20000, with \verb|max_weight| 25 being sufficient for almost every case. With default parameters (\verb|max_weight| 100, \verb|sos_limit| 20000), all consequences of each equation can be proven in at most $1$ second per equation, with the exception of laws $\Eq{450}$ and $\Eq{650}$ (and their duals), which take somewhat longer\footnote{For details, see \url{https://leanprover.zulipchat.com/\#narrow/channel/458659-Equational/topic/Which.20implications.20are.20harder.20for.20ATPs/near/547089094}.}.

\emph{Vampire} does not allow the imposition of a user-defined maximum weight, although its default LRS algorithm already implements a dynamic weight limit. According to \cite{janota2025experimentalresultsvampireequational}, \emph{Vampire}~5.0 can prove all true implications from the ETP, and can prove 99.97\% of them using less than 500 instructions per proof.

\subsection{Mace4}\label{Mace4}

The configuration of \emph{Mace4} likewise has a substantial impact on its running time, with different parameter settings resulting in differences of orders of magnitude. In particular, model-computation time is greatly affected by the \verb|selection_measure| and \verb|skolems_last| parameters. The \verb|selection_order| parameter is also relevant, but our empirical data indicate that \verb|selection_order| 2 is by far the best general choice. It is outperformed by \verb|selection_order| 0 only in a few cases, and even then the improvement in computing time is not that substantial. The effect of setting \verb|skolems_last| can range from slightly detrimental to highly advantageous: as an example, a model of $\Eq{272260}\nmodels\Eq{42323216}$ of size $7$ is found in $0.0$s with (2,4,Y), and in $4.0$s with (2,4,N). Here and in the following we use notation (a,b,S) to mean the configuration \verb|selection_order| a, \verb|selection_measure| b, and \verb|skolems_last| set (S$=$Y) or clear (S$=$N).

There is no universally optimal configuration (a,b,S): for each (potential) implication or theory, configurations may be ranked differently. However, in the ETP context, certain configurations have performed consistently better, roughly ordered as (2,3,Y), (2,1,Y), (2,4,Y), (2,3,N), (2,1,N), (2,4,N), (0,1,N), (0,1,Y), with the remaining configurations being largely avoidable. Problems involving several algebraic operations, or presenting some kind of symmetry (see \Cref{table:weakcentralgroupoids} on weak central groupoids), tend to perform better with (2,4,N/Y). Problems may also be considerably sensitive to changes in formulas that appear minor, for example adding $f^n(x)=x$ with different values of~$n$ can result in different optimal configurations.
\looseness=-1

\begin{table}
\centering
\caption{Largest size of $\Eq{1485}$ (weak central groupoids) models exhausted by each configuration in under 300s. The \texttt{skolems\_last} parameter turns out to be irrelevant for this problem.}
\label{table:weakcentralgroupoids}
\begin{tabular}{cccccc}
  \toprule
  Configuration &(2,4,Y/N) & (0,4,Y/N) & (2,3,Y/N) & (2,2,Y/N) & (2,1,Y/N) \\\midrule
  Size &14 (in \SI{127}{s}) & 14 (in \SI{166}{s}) & 13 (in \SI{290}{s}) & 9 (in \SI{33}{s}) & 7 (in \SI{8}{s}) \\
  \bottomrule
\end{tabular}
\end{table}

In addition, while in the ETP we typically searched for a single model to contradict a given implication, in some cases we sought all models of a certain size for some given theory, either to better understand a collection of base examples in order to extend them, or just to improve our understanding. Accordingly, it should be noted that a configuration that is fast for finding one model of a given size may be slower than others for exhausting the whole size. A notable example of this is (2,2,Y) for models of $\Eq{677}$ of size 9, which finds a single model in 160s, but is unable to exhaust the whole size after 20h (see \Cref{table:677models}).

When undertaking a long-running model search at a large size, it is advisable to determine in advance the potentially optimal configuration(s) by examining smaller sizes. Throughout the ETP, we have employed the following procedure. If time permits, exhaust the previous size with all configurations to make a choice. If that process is prohibitively time-consuming, then set the target number of desired configurations (e.g., for parallel runs), initialize the pool of configurations with all possibilities, and iterate:
\begin{enumerate}
\item Pick the next available smallest size.
\item Run all configurations on that size under a reasonable time limit and collect their running times, whether for finding a single model, a predetermined number of models, or for exhausting the size.
\item Based on the smallest time found, determine a statistically significant time threshold.
\item Remove from the pool those configurations whose time differs from the best one by more than the threshold (if any). End the loop once the number of configurations contained in the pool is the desired one.
\end{enumerate}
Be aware that this algorithm assumes that once a configuration outperforms another, this advantage carries over to larger sizes. This is not always the case (see \Cref{table:Schneider} for an example). Additionally, the algorithm may be further refined by considering combinations of the configurations with different subsets of formulas.

\begin{table}
\caption{Time to find one model of $\Eq{677}$ of size~$9$ and to exhaust the same size for different configurations. Some finite-setting formulas were included to speed up the search (see \Cref{finite setting}). The configurations not listed here were unable to find a model in \SI{4000}{s}.}
\label{table:677models}
\centering
\small
\sisetup{table-format=7.2, table-alignment-mode=format, table-number-alignment=center, table-align-text-post=false, table-align-text-pre=false}
\begin{tabular}{cSS}\toprule
Configuration & {One model} & {Exhaust size} \\\midrule
(2,1,N) & 16 \,\unit{s} & 190\,\unit{s} \\
(2,1,Y) & 16 \,\unit{s} & 210\,\unit{s}\\
(0,1,Y) & 20 \,\unit{s} & 149\,\unit{s}\\
(0,1,N) & 32 \,\unit{s} & 169\,\unit{s}\\
(2,2,Y) & 154 \,\unit{s} & {$>$}72000\,\unit{s} \\
(2,2,N) & 180 \,\unit{s} & {$>$}43000\,\unit{s} \\
(2,3,N) & 557 \,\unit{s} & 877\,\unit{s}\\
(2,3,Y) & 587 \,\unit{s} & 941\,\unit{s}\\
(2,4,Y) & 596 \,\unit{s} & 761\,\unit{s}\\
(2,4,N) & 620 \,\unit{s} & 730\,\unit{s}\\
(0,3,N) & 3746 \,\unit{s} & 4151\,\unit{s}\\
(0,3,Y) & 3777 \,\unit{s} & 4254\,\unit{s}\\
\bottomrule
\end{tabular}
\end{table}

\begin{table}
\caption{Comparison of the two best configurations for finding certain special $\Eq{677}$ models, both for exhausting each size and for finding a single model, with best times in boldface. Note that initially (2,4,N) seems the best configuration, but in the long run, (2,3,Y) outperforms it. The comparison is particularly misleading at size 16 when searching for a single model.  All times are in seconds.}
\label{table:Schneider}
\centering
\small
\sisetup{table-format=4.4, detect-weight=true, table-align-text-post=false}
\begin{tabular}{cS@{}SS@{}S}\toprule
& \multicolumn{2}{c}{Exhaustion} & \multicolumn{2}{c}{Finding one model} \\
\cmidrule(lr){2-3}\cmidrule(lr){4-5}
Size & {(2,3,Y)} & {(2,4,N)} & {(2,3,Y)} & {(2,4,N)}\\\midrule
9 & 0.14 & \bfseries 0.04 & 0.02 & \bfseries 0.0\\
10 & 0.32& \bfseries 0.12& {---} & {---}   \\
11 & 0.82& \bfseries 0.37& 0.15& \bfseries 0.01\\
12 & 2.17& \bfseries 1.2& {---} & {---}   \\
13 & 6.18& \bfseries 4.35& 2.74& \bfseries 0.91\\
14 & 15.47& \bfseries 14.39& {---} & {---}   \\
15 & \bfseries 38.82& 50.87& {---} & {---}   \\
16 & \bfseries 120& 175& 11.95& \bfseries 0.02\\
17 & \bfseries 306& 660& {---} & {---}   \\
18 & \bfseries 883& 2382& {---} & {---}  \\
19 & \bfseries 2300& {\!\llap{$>$}}3600 & \bfseries 846& 2082\\\bottomrule
\end{tabular}
\end{table}

\subsection{Finite Setting Case Study for Mace4}
The following example illustrates the finite-setting techniques of \Cref{finite setting} in a \emph{Mace4} search.
\begin{example}\label{atp-case-study}
We examine the model search for $\Eq{1518}\nmodels\Eq{47}$ using \emph{Mace4} (with $*$ standing for $\op$). It is now known that the smallest models have size~$15$. Initially, this search exceeded our ATPs expertise, hence we ended up constructing a theoretical model of size~$232$. If we restrict to the original formulas $\Eq{1518}$ and $\neg\Eq{47}$, the optimal configuration is (2,3,N), taking $2.5$ minutes to exhaust size~$10$ and already exceeding $7$~hours to exhaust size~$11$. Putting $S(x):=x\op x$, since $\Eq{1518}$ is $x = (y \op y) \op  (x \op  (y \op  x))$ we observe that $L_{S(y)}$ is surjective, hence injective. Accordingly we add \verb|(y*y)*x = (y*y)*z -> x = z| to \emph{Mace4} and then size 11 is exhausted in less than $5$ minutes. In addition, from the implication graph we observe that $\Eq{1518}$ implies $\Eq{359}$, which states $S(x) \op x = S(x)$; when added as \verb|(x*x)*x = x*x|, \emph{Mace4} exhausts size 11 in under 1s, with the optimal configuration shifting to (2,1,N). Size 12 is exhausted in 18s. We also note that if $S(x)=S(y)$ then $S(x)\op x = S(x) = S(y) = S(y)\op y$, yielding $x=y$ by injectivity of $L_{S(x)}$. Therefore $S$ is injective, we can add \verb|x*x = y*y -> x = y|, and exhaust size 12 in 13.5s. Also, in $\Eq{1518}$ we can substitute $x\mapsto S(y)\op x$ and cancel $S(y)$ on the left to get $\Eq{320858}$, with which we exhaust size 12 in under $2$s and size 13 in under $1$ minute. Moreover, as $S$ is surjective every element is a square, so the injectivity of $L_{S(x)}$ implies that of $L_x$, and we can add \verb|x*y = x*z -> y = z| to the search. But this actually worsens the performance! After some experimentation, we observe that we should either add the injectivity of $L_{S(x)}$ or that of $L_x$, but not both. It turns out that for size 13, injectivity of $L_{S(x)}$ performs slightly better, so we retain it. Additionally, the implication formula for this injectivity can be replaced with a new operation~\texttt{\textbackslash} by writing \texttt{x\textbackslash((x*x)*y) = y}. This further reduces the exhausting time for size 13 to $33$s. Now size 14 is exhausted in 30 minutes, and size 15 is well within reach.
\end{example}

\subsection{Discriminators for Isofilter}
\emph{Mace4}'s output may be redundant, in that many of the models found in a search may (and often will) be isomorphic to each other. To solve this problem, \emph{Mace4}'s output can be fed to another tool called \emph{Isofilter}, which returns a representative model of each isomorphism class. \emph{Isofilter} compares permutations of the models, restricted according to some discriminators: by default the only one used is the frequency with which each domain element appears in the operation tables\footnote{This discriminator is entirely ineffective with models having same number of copies of each element in their tables, such as quasigroups.}, but the user can include any number of other discriminators. \emph{Isofilter} ``as is'' handles models of size up to around 10 with ease; but the combinatorial explosion soon makes computations unfeasible for larger sizes, unless discriminators well suited to the problem are chosen. For example, \emph{Mace4} provides 10 models of size 15 for $\Eq{1518} \not\models \Eq{47}$, of which $6$ belong to different isomorphism classes. Vanilla \emph{Isofilter} requires 10 days and processes $3\cdot10^{13}$ permutations in order to determine these classes, while adding the discriminators \verb|x*x=x. (x*y)*z = x*(y*z). x*y=y. x*y=y*x.| reduces the task to $10^5$ permutations, completed in 0.02s.
